(1) 一个量子比特预选为 \(|i\rangle=c|0\rangle+s|1\rangle\)，后选择为 \(|f\rangle=c|0\rangle-s|1\rangle\)，其中 \(c=\cos(\theta/2),s=\sin(\theta/2),0<\theta<\pi/2\)。测量算符 \(A=\sigma_z\)，其本征值为 \(\pm1\)。求弱值 \(A_w=\langle f|A|i\rangle/\langle f|i\rangle\) 和零耦合时的后选择成功概率。
(2) 指针初态在位置 \(q\) 表象为归一化实Gaussian \(\psi_0(q)\propto e^{-q^2/(4\sigma^2)}\)，位置方差 \(\sigma^2\)，耦合 \(U=\exp(-igA\otimes p/\hbar)\)。在 \(g/\sigma\ll|\cos\theta|\) 下推导后选择后指针位置平均的一阶位移。
(3) 不采用弱耦合展开，写出未归一化后选择指针波函数，精确计算成功概率 \(P_{\rm succ}\) 与条件平均 \(\langle q\rangle_f\)，解释 \(\theta\to\pi/2\) 时弱值形式发散为何不意味着真实指针位移无限大。
(4) 在严格弱测量近似下计算成功率与位移放大平方的乘积 \(P_{\rm succ}A_w^2\)，结合Gaussian指针的量子Fisher信息说明后选择为什么不能凭空提高所有实验资源下的估计精度。
